
\section*{Future Perspectives}
\addcontentsline{toc}{section}{Future Perspectives}

The convergence of MoE architectures with healthcare AI is still in its early stages, and several transformative directions are poised to reshape both the technology and its clinical applications in the coming years.

\subsection*{Personalized expert routing for precision medicine}

Current MoE architectures route inputs based on data features, but future systems may incorporate patient-specific routing that adapts to individual patient characteristics over time. In precision medicine, the optimal analytical pathway for a patient depends not only on their current clinical state but also on their longitudinal history, genetic profile, and treatment response. A personalized routing mechanism could learn to adapt expert selection based on patient-specific factors, effectively creating a customized diagnostic or treatment recommendation pathway for each patient.

This vision extends the concept of personalized medicine from the biological level (matching treatment to molecular profile) to the computational level (matching analytical methodology to patient characteristics). For example, in cancer treatment planning, patients with different molecular subtypes might be routed through different expert pathways: one specializing in immunotherapy-responsive profiles, another in targeted therapy-responsive profiles, with the routing decision informed by the patient's genomic data, imaging features, and treatment history.

\subsection*{Federated MoE for multi-institutional collaboration}

The privacy-preserving training of MoE models across multiple healthcare institutions represents a compelling direction for overcoming the data silo problem in clinical AI. In a federated MoE framework, each institution could maintain local experts trained on its proprietary data, while shared experts capture cross-institutional patterns. The gating network would learn to route patient data to the most appropriate combination of local and shared experts, enabling personalized predictions that benefit from both institutional specificity and broad clinical knowledge.

This federated approach aligns with the collaborative nature of medical research, where multi-center studies are the gold standard for evidence generation. Federated MoE could enable continuous model improvement as new institutions join the network, with each institution contributing specialized expertise without compromising patient privacy. The technical challenges include communication-efficient training protocols, heterogeneous data distribution handling, and ensuring consistent expert behavior across institutions.

\subsection*{MoE in clinical foundation models}

The development of clinical foundation models, large-scale models pre-trained on diverse healthcare data that can be adapted to specific clinical tasks, will increasingly incorporate MoE architectures. Medical data spans multiple modalities (text, images, genomics, signals, structured records) and multiple scales (molecular, cellular, tissue, organ, patient, population), and MoE provides a natural framework for handling this diversity. Future clinical foundation models may employ hierarchical MoE structures, where higher-level experts handle modality routing and lower-level experts specialize within modalities.

The integration of MoE into clinical foundation models also enables efficient adaptation to new clinical settings. When a foundation model is deployed at a new institution, institution-specific experts can be trained on local data while the shared expert backbone remains fixed, enabling rapid adaptation with minimal local data requirements. This approach could dramatically accelerate the deployment of clinical AI across diverse healthcare settings worldwide.

\subsection*{Dynamic and adaptive expert creation}

Current MoE architectures employ a fixed number of experts determined at training time. Future systems may support dynamic expert creation, where new experts are spawned when the model encounters novel data patterns that are not well-served by existing experts. In healthcare, this capability could enable continuous learning as new diseases emerge, new treatments become available, or new patient populations are encountered. The gating network would evolve to accommodate new experts while maintaining the performance of existing specializations.

\subsection*{Integration with causal inference}

The combination of MoE architectures with causal inference methods represents an exciting frontier for clinical AI. Expert specialization in MoE models may capture distinct causal mechanisms underlying different disease subtypes or treatment responses. By incorporating causal structure into the gating mechanism (e.g., routing based on known disease mechanisms rather than purely statistical patterns), MoE models could provide not only predictions but also mechanistic explanations that align with clinical reasoning.

\subsection*{Regulatory science for modular AI}

The regulatory frameworks for clinical AI must evolve to accommodate the modular, conditional nature of MoE architectures. This includes developing standards for validating individual experts, characterizing routing behavior, and ensuring consistent performance across all input subpopulations. Regulatory science research in this area should establish evidence standards for MoE-based clinical AI that balance the need for comprehensive validation with the practical constraints of regulatory review.

\subsection*{Edge deployment and real-time inference}

Advances in model compression, expert pruning, and hardware-aware MoE design will enable the deployment of MoE-based clinical AI on edge devices at the point of care. In resource-limited settings, lightweight MoE models could provide clinical decision support on portable devices, with the sparse activation pattern of MoE naturally reducing computational requirements. Real-time applications, such as intraoperative guidance, ICU monitoring, and emergency triage, will benefit from the efficient inference enabled by MoE architectures.
